\section{Scope and boundary examples}
\label{sec:scope}
The center of the paper is a general sufficient-class acquired-geometry theorem, not a classification of all positive causal experiments. Its exact minimax assertion uses a fixed admissible witness with a common productive task law and an all-strategy cost identity, together with one common task-balancing program. It does not require a common primary row for all models or all strategies. If two histories in the same alleged driver state have success probabilities $1/4$ and $3/4$, a single exact driver row cannot represent both. The present upper still works when a proved conditional attempt bound and the task-measure balance remain valid; otherwise neither the finite-driver formula nor the new theorem is licensed.

Pointwise calibration is not necessary for these minimax laws. Static hidden signs have zero residual variance after fully revealing histories, and nearly pure quantum states can do the same. Positive current-report marks then fail on those fibers. The raw reference-plus-balance argument avoids that lower-variance premise; it does not assert that every such failed calibration can be repaired. Conversely, report-only balance and an unaffordable reference preparation are insufficient, as is a reference model whose supposedly unpredictable productive state is already leaked by a preceding report.

Action aliasing is harmless only when it preserves physical ports, costs, observable inputs and legal continuation kernels. Duplicating one physical command into two names does not increase Proposition~\ref{prop:physical}'s count. A phase-dependent dispatcher that can see an otherwise unavailable port state is an extra resource, not a relabeling. The proposition is not a universal minimal-controller characterization.

The finite-state upper is an explicit recursion: at the two physical port labels, issue the declared command; at the productive call, compute the virtual mark and finite decision; retain only the next label. It has no history tape, accurate hidden posterior, free real register, free time counter or externally randomized persistent program. Exact Borel descriptions and finite bit descriptions are distinguished by Proposition~\ref{prop:finite}. Optimal joint program/time/workspace complexity remains unproved.

The two intervention realizations score the post-command audit. The Gaussian theorem preserves the original current-sign target and uses no intervention on it. These are different declared tasks and are never exchanged by an untyped reduction. The reference-law geometry is derived from actual preparation and its cost is included, but its existence remains an experiment-class assumption. General learning of an unknown gate lower bound, unrestricted kernel learning, arbitrary unbounded losses, nonterminating protocols, infinite-dimensional unbounded operators, and large-deviation equivalence are not conclusions of the paper. The stopped results proved here use nonnegative costs and explicit conditional hazards; they are not substitutions of a random time into a fixed-horizon theorem.

The retained predictable-geometry manuscript contains the earlier pointwise calibration and conditional-envelope route, including nonexact acquired-mass bounds. It remains valid under its stated assumptions and is preserved in full as an optional companion. Earlier curvature, continuation, adaptive refinement, causal morphism and independent realization papers likewise keep their own hypotheses. None is used backwards to establish the central theorem here, and no historical archive is silently recertified by the present proof.
